\chapter{Avaliação por painel docente}
\label{cap:avaliacao}

% =====================================================================
% PAPEL DO CAPÍTULO: avaliação empírica das questões por professores.
% As Seções 6.1 a 6.6 relatam a primeira rodada (curadoria do acervo,
% cinco professores, sem CEP) como AVALIAÇÃO FORMATIVA PRELIMINAR:
% descritiva, sem inferência. O painel formal (Rodada 2: itens comuns,
% desenho cego, concordância entre avaliadores) está desenhado em
% plano_experimento_painel.md e figura na Seção 6.7 como etapa seguinte
% e limitação. A Seção 6.8 discute os resultados.
% Dados e script: sistema/medicoes/acervo-consolidado/avaliacao/respostas/.
% O perfil dos avaliadores e os aspectos éticos vêm do rascunho da
% aluna no Overleaf (out/2026), reescritos e integrados aqui.
% Os critérios de qualidade têm casa em 3.4; a hipótese, em 2.5; o erro
% silencioso, em 1.4 --- aqui só se referenciam.
% =====================================================================

A avaliação por professores é o principal instrumento de validação externa do sistema proposto.
Por mais que o Verificador confira cálculos e o Crítico aplique uma rubrica, a validação
conceitual e didática de uma questão é atribuição do docente, como se argumentou na
Seção~\ref{sec:papel-mediador}, e só quem conhece uma turma pode dizer se a questão serve para
ela. Por essa razão o sistema foi posto diante de professores de Matemática em exercício na rede
pública.

A avaliação aqui relatada é \emph{formativa e preliminar}. Trata-se de uma primeira rodada com
cinco professores, desenhada para revelar o que precisava ser corrigido no sistema antes de um
painel formal e não para estimar a qualidade das questões com rigor estatístico. O painel formal,
com itens comuns a todos os avaliadores, desenho cego para a detecção de erros e medida de
concordância à luz dos critérios da Seção~\ref{sec:criterios}, é apresentado na
Seção~\ref{sec:painel-formal} como etapa seguinte e como limitação deste trabalho.

\section{Objetivos e desenho da primeira rodada}
\label{sec:rodada1}
\label{sub:rodada1-desenho}

A primeira rodada teve três objetivos complementares. O primeiro foi verificar se os professores
encontravam erros matemáticos nas questões, fossem eles de cálculo, de formulação ou de
inconsistência entre o problema e o gabarito. O segundo foi estimar a adequação pedagógica de cada
questão, isto é, se o nível de dificuldade, o contexto e o raciocínio exigido eram compatíveis com
os estudantes reais desses professores. O terceiro foi colher a percepção docente sobre o sistema
como ferramenta de apoio ao planejamento.

Participaram cinco professores de Matemática da rede pública estadual, quatro em atuação no Mato
Grosso e um na Bahia, identificados neste texto como avaliadores 1 a 5. A seleção foi intencional
\citep{bogdan1994investigacao}, pois buscou-se reunir professores com experiência efetiva em sala
de aula e capazes de emitir julgamento qualificado sobre clareza, nível e viabilidade de uso das
questões. Todos lecionam tanto no Ensino Fundamental quanto no Ensino Médio, de modo que o seu
julgamento sobre adequação de nível é informado pelo conhecimento de toda a progressão curricular
da Educação Básica. O grupo é heterogêneo em dois aspectos que importam para a leitura dos
resultados. Quanto ao tempo de docência, os participantes acumulam de 4 a 20 anos de atuação, o
que reúne professores em início de carreira e professores com longa vivência das dificuldades dos
estudantes. Quanto à formação, três cursam mestrado, um é licenciado em Matemática e um tem
especialização em Educação Matemática. A Tabela~\ref{quad:perfil-avaliadores} sintetiza esse perfil.

% CONFERIR: o rascunho da aluna dizia "entre 4 e 15 anos, média de 8"
% no texto e 5, 5, 4, 10 e 20 anos no quadro. O texto acima segue o quadro.
\begin{table}[htbp]
  \centering
  \caption{Perfil dos avaliadores da primeira rodada}
  \label{quad:perfil-avaliadores}
  \small
  \begin{tabular}{c>{\raggedright\arraybackslash}p{5.6cm}ccc}
    \toprule
    Avaliador & Formação & Estado & Docência & Atuação \\
    \midrule
    1 & Licenciatura em Matemática, mestrado em andamento & MT & 5 anos & EF e EM \\
    2 & Licenciatura em Matemática, mestrado em andamento & MT & 5 anos & EF e EM \\
    3 & Licenciatura em Matemática & BA & 4 anos & EF e EM \\
    4 & Licenciatura em Matemática, mestrado em andamento & MT & 10 anos & EF e EM \\
    5 & Licenciatura em Matemática, especialização em Educação Matemática & MT & 20 anos & EF e EM \\
    \bottomrule
  \end{tabular}

  \smallskip
  {\footnotesize Todos atuam na rede pública estadual. EF e EM indicam Ensino Fundamental e Ensino
  Médio. Elaborado pela autora (2026).}
\end{table}

O material avaliado foi o acervo de 85 questões aprovadas pelo ciclo completo
(Seção~\ref{sec:produto-banco}). Em vez de replicar o acervo inteiro para cada pessoa, o que
tomaria várias horas de cada avaliador, adotou-se um delineamento de rotação com núcleo comum,
inspirado nos desenhos de avaliação por juízes \citep{streiner2015health}. Cada avaliador recebeu
um \emph{bloco-âncora} de seis questões comum a todos e um lote próprio de 15 ou 16 questões, num
total de 21 ou 22 itens e cerca de 75 minutos de trabalho por pessoa. A distribuição respeitou a
estratificação por habilidade, formato e dificuldade, e a ordem dos itens foi embaralhada para que
a âncora não ficasse no início. Foram colhidas 109 respostas, que cobrem as 85 questões, de modo que
cada questão do bloco-âncora foi julgada cinco vezes e cada uma das demais uma vez.

Cada item trazia o enunciado, as alternativas quando houvesse, o gabarito proposto e a habilidade
da BNCC declarada. Três informações foram omitidas de propósito. A resolução foi retirada porque
denunciaria um eventual erro e mudaria a tarefa. A garantia obtida (Seção~\ref{sec:garantia})
também foi omitida, já que a sua influência sobre o julgamento só pode ser examinada se ela não
estiver impressa na página. E o parecer do Crítico ficou de fora porque a sua presença mediria a
concordância com o agente e não o julgamento próprio do professor. Pelo mesmo motivo não se
informou a dificuldade declarada na especificação, para não ancorar a resposta.

Para cada item o avaliador respondia a quatro perguntas. Indicava se a questão tinha erro
matemático, com as opções não, sim com indicação do local e não sei dizer, e o nível de dificuldade
para os seus alunos, entre fácil, média, difícil ou fora do alcance da turma. Em seguida registrava
a decisão de uso, que podia ser aceita, aceita com ajuste ou recusada, e a justificativa, obrigatória
sempre que a decisão não fosse a aceitação. A instrução avisava que a amostra tinha qualidade
variável e podia conter erros matemáticos. Ao final três perguntas abertas tratavam do uso do sistema
no planejamento, do que faltou e do que atrapalhou. O instrumento completo está no
Apêndice~\ref{ap:instrumento}.

A participação foi voluntária e formalizada por um Termo de Consentimento Livre e Esclarecido, que
informava os objetivos da pesquisa, a natureza do material avaliado, a inexistência de respostas
certas ou erradas e a garantia de anonimato. Nenhum dado de identificação dos participantes é
apresentado, e a numeração dos avaliadores não corresponde a nenhuma característica identificável.
O material não contém informações sensíveis, não envolve estudantes e não representa risco aos
participantes. Ainda assim esta rodada não foi submetida a Comitê de Ética em Pesquisa, e por isso os
seus dados são tratados como avaliação formativa, descritivos e sem pretensão de generalização. O
painel formal, por sua vez, deverá ser submetido ao comitê nos termos da Resolução CNS nº 510/2016
\citep{brasil2016resolucao510}, que trata da ética na pesquisa em Ciências Humanas e Sociais.
% CONFERIR: confirmar com a aluna que o TCLE foi de fato assinado pelos
% cinco participantes da primeira rodada (o rascunho dela afirma que sim).

\section{Decisão de uso}
\label{sub:rodada1-decisao}

Das 109 respostas, 104 trazem decisão de uso. Seis em cada dez aceitam a questão como está, e oito
em cada dez a aceitam com ou sem ajuste, como mostra a Tabela~\ref{tab:rodada1-decisao}.

\begin{table}[htbp]
  \centering
  \caption{Decisão de uso, por avaliador}
  \label{tab:rodada1-decisao}
  \begin{tabular}{lrrrrr}
    \toprule
    Avaliador & Aceita & Com ajuste & Recusada & Em branco & Total \\
    \midrule
    1 & 18 & 1 & 0 & 3 & 22 \\
    2 & 9 & 7 & 6 & 0 & 22 \\
    3 & 8 & 7 & 6 & 1 & 22 \\
    4 & 18 & 3 & 1 & 0 & 22 \\
    5 & 9 & 5 & 6 & 1 & 21 \\
    \midrule
    Total & 62 & 23 & 19 & 5 & 109 \\
    \% das decisões & 59,6 & 22,1 & 18,3 & --- & \\
    \bottomrule
  \end{tabular}
\end{table}

Os avaliadores se dividem em dois perfis. Os avaliadores 1 e 4 aceitaram sem ajuste 36 das 41
questões que decidiram e recusaram apenas uma, enquanto os avaliadores 2, 3 e 5 aceitaram sem
ajuste 26 de 63 e recusaram 18. A diferença não parece ser de rigor matemático, já que o
avaliador~1 resolveu boa parte das questões à mão na própria folha. Ela parece ser de referência de
turma, como mostra a leitura das justificativas.

As recusas têm um traço comum. Todas as 19 recaem sobre questões que o próprio avaliador
classificou como \emph{difícil}, em 6 casos, ou \emph{fora do alcance da turma}, em 13, e todas as
justificativas invocam o nível de conhecimento dos alunos ou a distância entre a questão e a
realidade da sala de aula. Nenhuma recusa foi motivada por erro matemático. Das 16 respostas que
marcaram \emph{fora do alcance}, 13 terminaram em recusa e 3 em aceitação com ajuste, e das 25
respostas que marcaram \emph{fácil}, 23 foram aceitas sem ajuste. Na leitura dos professores,
portanto, a decisão de uso é sobretudo uma decisão sobre dificuldade.

O formato acompanha essa leitura. As questões discursivas foram recusadas em 13 de 48 decisões, ou
27\%, contra 6 de 56 nas de múltipla escolha, ou 11\%. As justificativas apontam enunciados longos,
com vários comandos encadeados, como o que mais pesa, e um dos avaliadores resumiu a impressão ao
escrever que ``os enunciados são estruturados para quem tem pleno domínio do conhecimento
matemático'' (avaliador~2).

Quanto à garantia obtida, as respostas sobre questões \emph{conferidas em parte} são poucas para
sustentar comparação. Foram 12, uma por questão, com cinco aceitas, cinco aceitas com ajuste, uma
recusada e uma em branco, e quatro dos cinco ajustes vieram do mesmo avaliador. O registro fica como
dado para o painel formal, desenhado para examinar essa diferença com estratos equilibrados.

\label{sub:rodada1-ancora}
As seis questões do bloco-âncora são as únicas em que os julgamentos podem ser comparados entre si,
e a Tabela~\ref{tab:rodada1-ancora} as reúne.

\begin{table}[htbp]
  \centering
  \caption{Decisões sobre as questões do bloco-âncora (A = aceita, J = aceita com ajuste,
  R = recusada)}
  \label{tab:rodada1-ancora}
  \begin{tabular}{llllccccc}
    \toprule
    & & & & \multicolumn{5}{c}{Avaliador} \\
    \cmidrule(l){5-9}
    Questão & Habilidade & Formato & Declarada & 1 & 2 & 3 & 4 & 5 \\
    \midrule
    Q007 & EM13MAT302 & múltipla escolha & fácil & A & A & A & A & J \\
    Q015 & EM13MAT403 & discursiva & média & A & R & R & A & J \\
    Q024 & EM13MAT501 & múltipla escolha & difícil & A & A & A & A & A \\
    Q034 & EM13MAT401 & múltipla escolha & fácil & A & A & J & A & J \\
    Q065 & EM13MAT306 & discursiva & difícil & A & R & R & J & R \\
    Q087 & EM13MAT507 & múltipla escolha & média & A & A & A & A & A \\
    \bottomrule
  \end{tabular}
\end{table}

Os extremos são nítidos. As duas questões unânimes são de múltipla escolha, e a Q024, declarada
difícil e de nível \emph{criar}, foi percebida como fácil ou média pelos cinco avaliadores. A Q065,
uma modelagem de maré por função cosseno com dois itens encadeados, foi recusada por três dos cinco,
sempre pelo nível da turma, embora a sua matemática esteja correta. Já a Q015, que pede a inversa de
uma exponencial e o seu domínio, divide o grupo com duas aceitações, um ajuste e duas recusas. As
decisões de todos os avaliadores sobre todas as questões estão no Apêndice~\ref{ap:instrumento}.

\section{Dificuldade declarada e dificuldade percebida}
\label{sub:rodada1-dificuldade}

O sistema classifica cada questão segundo a dificuldade pedida na especificação, e a
Tabela~\ref{tab:rodada1-dificuldade} cruza essa classificação com a que os professores deram
pensando nas suas turmas.

\begin{table}[htbp]
  \centering
  \caption{Dificuldade declarada pela especificação $\times$ dificuldade percebida pelos
  avaliadores (número de respostas)}
  \label{tab:rodada1-dificuldade}
  \begin{tabular}{lrrrrr}
    \toprule
    & \multicolumn{5}{c}{Percebida} \\
    \cmidrule(l){2-6}
    Declarada & Fácil & Média & Difícil & Fora do alcance & Em branco \\
    \midrule
    Fácil & 13 & 10 & 6 & 3 & 1 \\
    Média & 6 & 16 & 9 & 6 & 1 \\
    Difícil & 6 & 14 & 11 & 7 & 0 \\
    \bottomrule
  \end{tabular}
\end{table}

Das 107 respostas com dificuldade marcada, 40 coincidem com a declarada, 41 a colocam acima e 26
abaixo, ou seja, 37\%, 38\% e 24\% respectivamente. O desvio predominante é para cima. Quatro dos
cinco avaliadores julgaram mais questões acima do que abaixo do nível declarado, os avaliadores 2 e
5 com folga, por 13 contra 4 e 13 contra 2, e os avaliadores 1 e 3 por margem pequena, por 4 contra
2 e 8 contra 7. Só o avaliador~4 deslocou a escala para baixo, por 3 contra 11. O efeito é mais
visível no degrau de baixo, pois das 32 questões declaradas fáceis que receberam resposta apenas 13
foram percebidas assim e 9 foram julgadas difíceis ou fora do alcance.

As justificativas indicam a origem do desvio. Cada professor julga pela própria turma, e as turmas
estão aquém do que a habilidade supõe, como registraram dois deles ao escrever que ``meus alunos não
sabem nem trigonometria no triângulo retângulo'' (avaliador~2) e que ``no Ensino Médio, são três
aulas de Matemática por semana'' (avaliador~5). A escala do Gerador é a de um estudante que cumpre a
habilidade da BNCC no ano previsto, e o seu degrau mais baixo, o fácil, definido como aplicação
direta de um conceito em poucos passos, ainda pressupõe esse estudante.

Em resposta o sistema ganhou um quarto nível, o \emph{básico}, abaixo do fácil. Ele pede um único
passo com aplicação imediata de uma definição ou fórmula, números pequenos, enunciado curto em
linguagem cotidiana e uma só pergunta, sem itens encadeados, e incorpora assim também as queixas de
linguagem registradas na Seção~\ref{sub:rodada1-comentarios}. Antes de ser oferecido ao professor
o nível foi posto à prova com provedor real em um desenho pareado. Para cada uma das quinze
habilidades a especificação fácil do acervo foi gerada duas vezes na mesma execução, uma como fácil
e outra como básica, com todo o resto igual. A Tabela~\ref{tab:basico} resume o resultado.

\begin{table}[htbp]
  \centering
  \caption{Questões fáceis e básicas geradas a partir das mesmas 15 especificações
  (indicadores de texto; medianas sobre os 14 pares com as duas questões aprovadas)}
  \label{tab:basico}
  \begin{tabular}{lrr}
    \toprule
    & Fácil & Básico \\
    \midrule
    Aprovadas pelo ciclo & 15 de 15 & 14 de 15 \\
    Palavras no enunciado & 134,5 & 75,5 \\
    Pares em que o básico é mais curto & \multicolumn{2}{c}{13 de 14} \\
    Discursivas em que o básico eliminou os itens encadeados & \multicolumn{2}{c}{5 de 7} \\
    Pares com menos números no enunciado básico & \multicolumn{2}{c}{9 de 14} \\
    Nota $\le 3$ em alinhamento à BNCC (Crítico) & 3 de 14 & 3 de 14 \\
    \bottomrule
  \end{tabular}
\end{table}

O nível básico encurta o enunciado quase sempre, desfaz a cadeia de itens na maior parte das
discursivas e nunca faz o contrário, e segundo o Crítico a simplificação não custou alinhamento à
habilidade. O limite apareceu onde era de esperar. Nas três especificações de nível cognitivo
\emph{criar} só uma simplificou, enquanto outra foi descartada, com o Crítico apontando nas três
tentativas que um item de um passo não exige criação, e a terceira saiu básica apenas no nome. Nos
níveis \emph{entender}, \emph{aplicar} e \emph{analisar} as doze especificações foram aprovadas. Por
essa razão a interface passou a avisar o professor quando ele combina o nível básico com o processo
\emph{criar}, sem impedi-lo de fazê-lo. Os indicadores, contudo, são de texto e não de dificuldade.
Eles mostram que a questão básica é mais curta, menos encadeada e tem menos números, mas não que a
turma a ache fácil, e essa é uma pergunta para o painel formal.

\section{Apontamentos de erro matemático}
\label{sub:rodada1-erros}

Em 99 das 109 respostas o avaliador declarou que a questão não tinha erro, em 6 deixou a pergunta
em branco e em 4 marcou que tinha. Nenhum dos quatro apontamentos se confirmou.

\begin{itemize}
  \item Na Q074, que associa uma progressão aritmética a uma função afim, o avaliador~2 indicou
  erro ``no gabarito'' e marcou o mês 215 em vez do 216. O gabarito está certo, pois com
  $f(n) = 340 + 45(n-1)$ tem-se $f(215) = 9\,970$ e $f(216) = 10\,015$. O mês 215 é justamente o
  distrator previsto para quem resolve $45(n-1) > 9\,660$ e não converte $n-1$ em $n$.
  \item Nas Q073 e Q084 o avaliador~5 marcou erro sem indicar onde. As duas estão corretas, e em
  ambas o avaliador anotou ao lado uma redação da habilidade declarada diferente da do catálogo, de
  modo que o apontamento parece referir-se a essa redação e não à questão.
  \item Na Q023 o avaliador~4 marcou erro sem indicar onde e ainda assim aceitou a questão sem
  ajuste. A questão está correta, e como em todas as outras 21 ele marcou a caixa vizinha, que
  corresponde à resposta \emph{não}, é possível que a marca tenha caído na caixa errada.
\end{itemize}

O dado mais importante desta seção é o inverso, uma questão com erro que ninguém apontou. Na Q061
uma cultura de bactérias começa com 200 indivíduos e dobra a cada hora, e o enunciado define $a_n$
como a quantidade \emph{ao final} da $n$-ésima hora. Logo $a_1 = 400$, $a_n = 200 \cdot 2^n$ e
$a_6 = 12\,800$. O gabarito, porém, afirma $a_1 = 200$, $f(n) = 200 \cdot 2^{n-1}$ e
$a_6 = 6\,400$, e está portanto deslocado em um termo. O avaliador~1 sublinhou a definição de $a_n$
e escreveu na margem ``$a_1 = 400$, $a_2 = 800$'' e ``$f(0) = 200$'', que é a conta certa, e mesmo
assim marcou \emph{não} na pergunta sobre erro e aceitou a questão.

O caminho desse erro pelo sistema é instrutivo. O Gerador formalizou a progressão com $a_1 = 200$ e
razão $2$, o Verificador conferiu as quatro afirmações dessa formalização e todas estavam
consistentes entre si, e o Crítico deu nota 5 à clareza. A questão saiu com a garantia
\emph{gabarito conferido}. O Verificador não errou no que lhe cabia, pois ele confere a formalização
declarada e não o enunciado (Seção~\ref{sec:ag-verificador}). O erro está na tradução da prosa para a
formalização, etapa que nenhum dos agentes confere e que o professor, mesmo fazendo a conta certa,
não converteu em apontamento. Trata-se de um erro silencioso no sentido da
Seção~\ref{sec:erro-silencioso}, e o caso confirma na prática o alcance que a
Seção~\ref{sec:garantia} atribui à garantia \emph{gabarito conferido}, que é a consistência do
gabarito com a formalização declarada e não com o enunciado. A demonstração é direta, pois com o
gabarito corrigido e a formalização ajustada a ele, com $a_1 = 400$, o Verificador também aprova a
questão. As duas respostas, incompatíveis entre si, passam pela mesma conferência para o mesmo
enunciado.

Com um único caso não há taxa a relatar, já que a rodada não continha erros postos de propósito e o
acervo não foi conferido à mão por inteiro. O que o caso estabelece é a existência do ponto cego e
não a sua frequência.

\section{Comentários e questionário final}
\label{sub:rodada1-comentarios}

Os comentários livres convergem em cinco temas.

O tema dominante é a \textbf{distância entre a questão e a turma}, presente em todas as recusas e
nas respostas de três avaliadores ao questionário final. Para um deles as questões ``estão adequadas;
porém, não são compatíveis com a realidade e com o nível de conhecimento dos estudantes''
(avaliador~5), e outro afirmou que usaria o sistema ``com um nivelamento de dificuldade das
questões'' (avaliador~3).

Em seguida aparecem a \textbf{linguagem e a estrutura do enunciado}. Os avaliadores apontaram
enunciados longos e técnicos, com comandos encadeados, e pediram ``uma linguagem menos complexa''
(avaliador~3) ou que se adequasse ``a linguagem ao público'' (avaliador~2, na Q055), enquanto outro
registrou ``texto com partes confusas'' (avaliador~4, na Q068).

O terceiro tema é a \textbf{ausência de figuras e de resolução}. Três dos cinco avaliadores pediram
representação visual sem terem combinado entre si, falando em ``questões com o esboço do gráfico''
(avaliador~1), em ``imagens ou uma descrição bem detalhada dessas'' (avaliador~2) e em ``figuras
matemáticas que ilustrassem o enunciado'' (avaliador~5). O sistema hoje produz apenas texto, e o
pedido é mais agudo justamente nas habilidades que tratam de conversão entre os registros algébrico
e gráfico. O avaliador~4 pediu ainda ``uma seção de resolução da questão além do gabarito proposto'',
que o sistema produz mas que foi retirada deste material por desenho.

O quarto tema é a \textbf{correspondência com a BNCC}. Dois avaliadores questionaram as habilidades
declaradas. O avaliador~4 escreveu que ``alguns códigos de habilidade não [batem] com [os] da BNCC''
e destacou três deles, \texttt{EM13MAT403}, \texttt{EM13MAT404} e \texttt{EM13MAT405}, e o
avaliador~5 reescreveu ao lado de duas questões uma redação diferente para \texttt{EM13MAT306} e
\texttt{EM13MAT503}. Os códigos e textos do catálogo conferem palavra por palavra com o documento da
BNCC que lhe serve de fonte (Seção~\ref{sec:bncc}). A divergência sugere que os professores trabalham
com outra referência, seja outra edição da Base, seja um documento curricular derivado dela, e o
ponto merece ser esclarecido antes do painel formal, porque a habilidade declarada é a âncora do
critério de alinhamento curricular.

Por fim surgiram observações sobre os \textbf{padrões do próprio gerador}. O avaliador~1 notou que
``a resposta [está] sempre na letra A'' e que a Q087 repete a progressão da Q045, e os dois
apontamentos se confirmam no acervo. Em 38 das 43 questões de múltipla escolha aprovadas a
alternativa correta era a primeira, e as duas questões usam a mesma progressão, com $a_1 = 7$ e
$r = 4$. A repetição alcança também os contextos, pois culturas de bactérias aparecem em 12 das 85
questões e corridas por aplicativo em 7, o que o avaliador~2 anotou como sugestão ao escrever
``insira no prompt o comando para diversificar e variabilizar o contexto''.

Sobre o uso do sistema no planejamento, quatro avaliadores responderam que o usariam, com
ressalvas. Um deles destacou a ``praticidade na elaboração de atividades e provas'' (avaliador~1),
outro disse que ``usaria, no entanto a revisão humana teria que ser criteriosa'' (avaliador~2), um
terceiro que ``usaria com um nivelamento de dificuldade'' (avaliador~3) e o quarto que o sistema
``facilita no planejamento das aulas'' (avaliador~4). O avaliador~5 não respondeu à pergunta de
forma interpretável.

\section{Limitações e consequências da primeira rodada}
\label{sub:rodada1-limitacoes}

Cinco avaliadores, cada um com um lote diferente, não sustentam medida de concordância nem
comparação entre estratos, e os números desta rodada descrevem o que aconteceu nela sem estimar
propriedades do sistema. O bloco-âncora, único conjunto comum, tem apenas seis questões. O desenho
também não mede detecção de erro, porque a instrução avisava que podia haver erros mas a amostra não
continha erros postos de propósito, e o único erro real foi encontrado depois, na análise das
respostas.

Ainda assim a rodada produziu consequências concretas para o sistema, todas verificáveis no acervo.
Quatro delas foram incorporadas. A ordem das alternativas passou a ser sorteada depois da geração,
com as referências por letra no gabarito e na resolução renumeradas junto, e na medição da
Tabela~\ref{tab:basico} a alternativa correta caiu nas posições a, b, c e d em 5, 5, 3 e 2 das 15
questões de múltipla escolha, contra 38 de 43 na primeira posição antes da mudança. O nível
\emph{básico} foi acrescentado e medido. A Q061 foi corrigida no acervo, preservado o registro
original do erro. E o alcance da garantia \emph{gabarito conferido} passou a ser declarado com
precisão na Seção~\ref{sec:garantia}. Dois pontos ficam em aberto, que são diversificar contextos e
dados numéricos entre ciclos e esclarecer a referência curricular com que os professores conferem as
habilidades.

\section{O painel formal}
\label{sec:painel-formal}

A primeira rodada não cumpre sozinha o objetivo de avaliar as questões por painel com instrumento
estruturado, e tampouco testa a hipótese arquitetural da Seção~\ref{sec:hipotese}. Para isso foi
desenhado um painel formal, cujo protocolo foi fixado antes de qualquer coleta, de modo que o plano
de análise não possa ser ajustado aos dados. O painel não foi aplicado no prazo desta dissertação, e
por isso é apresentado aqui como etapa seguinte e registrado como limitação.

O painel se organiza em torno de três perguntas. A primeira é se as questões servem, isto é, qual é a
qualidade percebida pelos professores em cada critério da rubrica e com que frequência eles as
aceitariam para uso real. A segunda, que é a que testa a hipótese, é se a verificação simbólica
intercepta erros que o professor não percebe. A terceira é se a garantia declarada significa alguma
coisa, ou seja, se questões marcadas como \emph{gabarito conferido} são julgadas de modo diferente
das marcadas como \emph{conferido em parte}. Se não houver diferença o rótulo é decorativo, e é melhor
saber disso.

O projeto original previa um desenho cego que misturasse questões do sistema com itens de bancos
humanos, como os do ENEM e da OBMEP. Esse contraste foi substituído por outro, interno ao sistema,
por duas razões. Itens do ENEM passam por meses de edição profissional e pré-testagem, e perder para
eles não diria nada sobre a arquitetura. Além disso, a tese desta dissertação não é a de que o
sistema gera questões tão boas quanto as do ENEM, e sim a de que a verificação simbólica intercepta um
erro que o modelo produziria e o professor não notaria. O desenho cego passa então a contrastar
questões aprovadas pelo ciclo completo com questões que o Verificador reprovou, produzidas pelo mesmo
gerador, no mesmo formato e sobre as mesmas habilidades, cada uma com o erro confirmado à mão.

O instrumento separa detectar erro de julgar qualidade, porque uma escala de 1 a 5 em correção
matemática não mede detecção. Um professor que não percebeu o erro marca 5 com sinceridade, e o dado
se torna indistinguível do de quem conferiu e aprovou. Por isso o material se divide em dois blocos
que todos os professores avaliam por inteiro, já que medidas de concordância exigem avaliadores
completos sobre os mesmos itens. O primeiro bloco, cego, tem 12 itens, metade deles aprovados e
metade com erro confirmado, e pergunta apenas se a questão tem erro matemático e onde. O segundo tem
24 itens estratificados por habilidade, formato, dificuldade e garantia obtida e pede notas de 1 a 5
nos cinco critérios didáticos da Seção~\ref{sec:criterios}, a decisão de uso e um comentário. Ao
final um questionário traça o perfil do participante e a sua percepção sobre o uso de inteligência
artificial na elaboração de questões, e só então se revela que as questões foram geradas por um
sistema com verificação simbólica. O tempo estimado é de cerca de 70 minutos, e o painel prevê de 8 a
15 professores.

O plano de análise acompanha as três perguntas. Para a qualidade, calculam-se média e desvio por
critério, a distribuição das decisões de uso e a concordância entre avaliadores. Para a detecção, a
medida central é a proporção de questões com erro em que o professor o apontou, contrastada com a taxa
de falso alarme nas questões corretas, o que separa detecção de desconfiança generalizada. Se essa
taxa for baixa, o erro silencioso deixa de ser uma premissa citada da literatura e passa a ser um
resultado medido com professores brasileiros e questões do Ensino Médio. Para a garantia, comparam-se
as notas e as decisões entre os estratos, tendo como hipótese nula a ausência de diferença.

Um obstáculo prático já é conhecido e precisa ser enfrentado antes da coleta. Como se viu na
Seção~\ref{sec:custo}, nas execuções com o modelo usado no acervo nenhuma reprovação do Verificador
correspondeu a um gabarito errado, de modo que os espécimes de erro não surgem espontaneamente em
quantidade suficiente. Será preciso produzi-los de propósito, guardando a primeira tentativa
reprovada em gerações feitas com a revisão desligada ou com modelos menores, e conferir cada um à mão,
porque um espécime cujo erro não existe destruiria o padrão de referência do bloco.

\section{Discussão}
\label{sec:discussao-avaliacao}

Lidos em conjunto, os resultados da primeira rodada permitem algumas conclusões, desde que se
respeitem os limites de uma avaliação formativa com cinco professores.

A primeira diz respeito à aceitação. Oito em cada dez decisões aceitaram a questão com ou sem
ajuste, e nenhuma recusa foi motivada por erro matemático. Por outro lado o Crítico havia aprovado
todas as 85 questões, e os professores recusaram 18\% das decisões. A distância entre as duas
avaliações não é de rigor e sim de referência, pois o Crítico julga a adequação ao nível tendo como
parâmetro o estudante que a BNCC descreve, enquanto o professor julga pela turma que tem diante de si.
Esse resultado reforça a posição assumida na Seção~\ref{sec:papel-mediador}, segundo a qual nenhum
agente automatizado substitui o conhecimento que o professor tem dos seus alunos. A resposta do
sistema foi oferecer ao professor um meio de declarar essa realidade na especificação, que é o nível
básico, em vez de pretender que o Crítico a adivinhe.

A segunda conclusão diz respeito ao erro silencioso. O único erro matemático do acervo atravessou
o Verificador, o Crítico e um professor que fez a conta certa na margem da folha. Esse episódio
confirma a premissa que motivou o trabalho, porque um gabarito errado redigido com fluência passou
por um leitor atento, e ao mesmo tempo delimita com precisão o que a verificação simbólica garante.
Ela elimina o cálculo errado sobre um problema bem formalizado, mas não protege contra a leitura
errada do enunciado no momento da formalização. Daí decorre uma consequência de projeto, já que
conferir a fidelidade da formalização ao enunciado é o próximo ponto a atacar, e uma consequência de
comunicação, que é declarar ao professor o alcance exato de cada rótulo, como passou a fazer a
Seção~\ref{sec:garantia}.

A terceira conclusão diz respeito ao papel de cada agente. Com o modelo usado no acervo, o ciclo de
revisão trabalhou sobretudo na dimensão didática, pois a maior parte das revisões foi pedida pelo
Crítico. A contribuição do Verificador apareceu menos como correção de erros do Gerador e mais como
garantia declarada, já que 72 das 85 questões chegaram ao professor com o gabarito recalculado de
forma independente. Avaliar se essa garantia muda o julgamento docente, e se a verificação torna
viáveis modelos menores que erram mais, são justamente as perguntas que o painel formal e a
comparação entre modelos foram desenhados para responder.

Por fim, a rodada mostrou o valor de submeter o sistema a professores antes de declará-lo pronto.
Pedidos como figuras, linguagem mais simples, diversidade de contextos e um nível abaixo do fácil não
apareceriam em nenhuma medição automática, e vários deles já se converteram em mudanças no sistema.
